%% kapa: exam=2 round=33 year=2022 period=1 subject=appraisal_practice subject_ko=감정평가실무 seq=3 number=3 points=20
%% source: Q-Net 감정평가사 기출문제 (cst00309) artlSeq=5210838 file=raw/1008/5210838/2022년 제33회 감정평가사 2차시험 1교시 문제지(감정평가실무).pdf
%% conversion: kordoc 4.24.51 -> tools/gen2.py (auto); review=proofread-gpt-astra; flags=none
\begin{longquestion}{【문제 3】}
감정평가사 甲은 토지의 장기임차권을 매입하여 지상에 공장 건물을 신축하여 사업체를 운영하고 있는 사업자 乙로부터 일반거래 목적의 감정평가를 의뢰받았다. 이해관계인은 공정한 자산 가액의 산정을 위하여 복수의 감정평가를 요구하고 있다. 제시된 자료를 활용하여 각 물음에 답하시오. (20점) \par
\medskip\noindent 물음 1) 대상 토지의 장기임차권 매입금액을 기준으로 한 시산가액을 산정하시오. (6점) \par
\medskip\noindent 물음 2) 토지의 장기임차권 거래사례 중 감정평가에 활용할 거래사례 하나를 선정하여 그 사유를 설명하고, 이를 기준으로 한 시산가액을 산정하시오. (6점) \par
\medskip\noindent 물음 3) 산정된 시산가액을 검토하여 감정평가액을 결정하시오. (4점) \par
\medskip\noindent 물음 4) 복수감정평가의 장단점을 설명하시오. (4점) \par
< 자료 1 > 기본적 사항 \par
1. 기준가치: 시장가치 \par
2. 기준시점: 2022년 7월 16일 \par
3. 대상 물건의 개황 \par
{\small\centering\renewcommand{\arraystretch}{1.25}
\resizebox{\ifdim\width>\linewidth\linewidth\else\width\fi}{!}{
\begin{tabular}{|c|c|c|c|c|c|c|c|}
\hline
\makecell{소재지\\지번} & \makecell{지목} & \makecell{면적(m\textsuperscript{2})} & \makecell{용도지역} & \makecell{이용상황} & \makecell{도로접면} & \makecell{형상\\지세} & \makecell{주위환경} \\
\hline
\makecell{A시 B동\\110} & \makecell{공장용지} & \makecell{2,000} & \makecell{일반공업지역} & \makecell{공업용} & \makecell{중로각지} & \makecell{가장형\\평지} & \makecell{일반공장\\지대} \\
\hline
\end{tabular}}\par}

4. 대상 토지의 장기임차권 내용 \par
1) 계약일: 2011년 1월 1일 \par
2) 계약기간: 50년 (2011년 1월 1일 ～ 2060년 12월 31일) \par
3) 매입금액: 120,000원/m\textsuperscript{2} \par
4) 계약내용: 토지의 장기임차권 매입금액은 계약일에 토지의 소유권자에게 일괄 지급하고, 계약기간 50년 동안 토지상에 건물과 공작물의 설치 등 토지를 안정적으로 사용할 수 있으며, 장기임차권 만료일이 경과되면 토지와 건물 등 모든 시설의 소유권은 토지의 소유권자에게 무상으로 반납된다. \par
< 자료 2 > 시점수정 자료 \par
1. 토지의 장기임차권에 대한 변동률이 고시되지 아니하여 당해 시의 공업지역 지가변동률을 적용함 \par
2. 지가변동률 \par
{\small\centering\renewcommand{\arraystretch}{1.25}
\resizebox{\ifdim\width>\linewidth\linewidth\else\width\fi}{!}{
\begin{tabular}{|c|c|c|}
\hline
\makecell{구분} & \makecell{A시 공업지역(\%)} & \makecell{비고} \\
\hline
\makecell{2011. 1. 1. ～ 2011. 12. 31.} & \makecell{5.001} & \makecell{누계} \\
\hline
\makecell{2012. 1. 1. ～ 2012. 12. 31.} & \makecell{6.505} & \makecell{누계} \\
\hline
\makecell{2013. 1. 1. ～ 2013. 12. 31.} & \makecell{6.312} & \makecell{누계} \\
\hline
\makecell{2014. 1. 1. ～ 2014. 12. 31.} & \makecell{7.322} & \makecell{누계} \\
\hline
\makecell{2015. 1. 1. ～ 2015. 12. 31.} & \makecell{8.457} & \makecell{누계} \\
\hline
\makecell{2016. 1. 1. ～ 2016. 12. 31.} & \makecell{5.023} & \makecell{누계} \\
\hline
\makecell{2017. 1. 1. ～ 2017. 12. 31.} & \makecell{4.505} & \makecell{누계} \\
\hline
\makecell{2018. 1. 1. ～ 2018. 12. 31.} & \makecell{3.255} & \makecell{누계} \\
\hline
\makecell{2019. 1. 1. ～ 2019. 12. 31.} & \makecell{2.975} & \makecell{누계} \\
\hline
\makecell{2020. 1. 1. ～ 2020. 12. 31.} & \makecell{2.523} & \makecell{누계} \\
\hline
\makecell{2021. 1. 1. ～ 2021. 12. 31.} & \makecell{2.350} & \makecell{누계} \\
\hline
\makecell{2022. 1. 1. ～ 2022. 05. 31.} & \makecell{1.244} & \makecell{누계} \\
\hline
\makecell{2022. 5. 1. ～ 2022. 05. 31.} & \makecell{0.198} & \makecell{5월} \\
\hline
\end{tabular}}\par}

- 지가변동률은 2022년 6월 이후는 고시되지 않아서 5월 지가변동률을 연장ㆍ추정하여 적용함 \par
< 자료 3 > 잔존가치율 산정 자료 \par
토지 장기임차권의 상각은 정액법에 따르고, 상각은 월단위 만월 상각을 적용하며, 잔존가치율 산정은 소수점 넷째 자리를 반올림하여 셋째 자리까지 표기한다. \par
< 자료 4 > 토지의 장기임차권 거래사례 \par
대상 토지의 인근지역에 위치하고 확인 가능한 토지의 장기임차권 거래사례는 다음 표와 같으며, 검토 결과 거래가액은 적정한 것으로 판단됨 \par
{\small\centering\renewcommand{\arraystretch}{1.25}
\resizebox{\ifdim\width>\linewidth\linewidth\else\width\fi}{!}{
\begin{tabular}{|c|c|c|c|c|c|c|c|c|c|c|}
\hline
\makecell{기호} & \makecell{소재지\\지번} & \makecell{지목} & \makecell{면적(m\textsuperscript{2})} & \makecell{용도지역} & \makecell{이용\\상황} & \makecell{도로접면} & \makecell{형상\\지세} & \makecell{단가\\(원/m\textsuperscript{2})} & \makecell{계약일} & \makecell{계약기간} \\
\hline
\makecell{가} & \makecell{A시\\B동\\115} & \makecell{공장용지} & \makecell{2,800} & \makecell{일반공업지역} & \makecell{공업용} & \makecell{광대한면} & \makecell{가장형\\평지} & \makecell{130,000} & \makecell{2011.07.16.} & \makecell{계약일로부터\\50년} \\
\hline
\makecell{나} & \makecell{A시\\B동\\210} & \makecell{공장용지} & \makecell{3,000} & \makecell{일반공업지역} & \makecell{공업용} & \makecell{중로\\한면} & \makecell{정방형\\평지} & \makecell{280,000} & \makecell{2022.07.01.} & \makecell{계약일로부터\\50년} \\
\hline
\makecell{다} & \makecell{A시\\B동\\220} & \makecell{공장용지} & \makecell{800} & \makecell{일반공업지역} & \makecell{상업용} & \makecell{중로\\한면} & \makecell{정방형\\평지} & \makecell{450,000} & \makecell{2022.07.01.} & \makecell{계약일로부터\\50년} \\
\hline
\makecell{라} & \makecell{A시\\B동\\230} & \makecell{공장용지} & \makecell{1,200} & \makecell{준공업지역} & \makecell{공업용} & \makecell{중로각지} & \makecell{세장형\\평지} & \makecell{360,000} & \makecell{2022.07.05.} & \makecell{계약일로부터\\50년} \\
\hline
\end{tabular}}\par}

< 자료 5 > 토지의 개별요인 \par
1. A시 B동 공업지대 가로조건, 획지조건의 개별요인 비교치는 아래와 같으며, 나머지 개별요인은 대등한 것으로 상정함 \par
2. A시 B동 공업지대의 가로조건 \par
{\small\centering\renewcommand{\arraystretch}{1.25}
\resizebox{\ifdim\width>\linewidth\linewidth\else\width\fi}{!}{
\begin{tabular}{|c|c|c|c|}
\hline
\makecell{구분} & \makecell{광대한면} & \makecell{중로각지} & \makecell{중로한면} \\
\hline
\makecell{광대한면} & \makecell{1.00} & \makecell{0.97} & \makecell{0.95} \\
\hline
\makecell{중로각지} & \makecell{1.03} & \makecell{1.00} & \makecell{0.98} \\
\hline
\makecell{중로한면} & \makecell{1.05} & \makecell{1.02} & \makecell{1.00} \\
\hline
\end{tabular}}\par}

3. A시 B동 공업지대의 획지조건 \par
{\small\centering\renewcommand{\arraystretch}{1.25}
\resizebox{\ifdim\width>\linewidth\linewidth\else\width\fi}{!}{
\begin{tabular}{|c|c|c|c|}
\hline
\makecell{구분} & \makecell{정방형} & \makecell{가장형} & \makecell{세장형} \\
\hline
\makecell{정방형} & \makecell{1.00} & \makecell{1.01} & \makecell{0.99} \\
\hline
\makecell{가장형} & \makecell{0.99} & \makecell{1.00} & \makecell{0.98} \\
\hline
\makecell{세장형} & \makecell{1.01} & \makecell{1.02} & \makecell{1.00} \\
\hline
\end{tabular}}\par}

< 자료 6 > 기타 사항 \par
1. 시점수정치인 지가변동률은 백분율로서 소수점 넷째 자리를 반올림하여 셋째 자리까지 표기한다. \par
2. 단가 산정은 천원 미만 단위에서 반올림하여 천원 단위까지 표기한다. \par
3. 개별요인 산정은 소수점 넷째 자리를 반올림하여 셋째 자리까지 표기한다. \par
4. 주어진 자료 이외의 사항은 고려하지 아니한다. \par
\end{longquestion}
